\subsection{Relevamiento exploratorio mediante encuesta}
\label{subsec:ml-encuesta-metodo}

Como complemento de las verificaciones técnicas se aplicó una encuesta exploratoria de carácter descriptivo, durante la etapa de documentación del proyecto, para conocer percepciones sobre la seguridad en marketplaces y el interés potencial en herramientas preventivas. Se empleó un formulario en línea con \textbf{once preguntas cerradas}: opciones únicas, escalas de 1 a 5 y selección múltiple. Las preguntas cubrieron hábitos de compra, criterios para valorar publicaciones, experiencias de desconfianza, percepción de los mecanismos actuales de protección, utilidad potencial y preocupaciones sobre herramientas independientes de evaluación.

La difusión se realizó mediante contactos del autor, por lo que la muestra es \textbf{no probabilística por conveniencia}. El corte considerado comprende \textbf{61 respuestas completas} recibidas entre el 10 y el 15 de noviembre de 2026. Se analizaron frecuencias absolutas y porcentajes sobre el total de participantes; las preguntas que admitían varias elecciones se contabilizaron por opción, sin sumar sus porcentajes como si fueran categorías mutuamente excluyentes.

El objetivo de esta actividad fue explorar la \textbf{necesidad percibida y el interés declarado} por la propuesta, no medir satisfacción con el prototipo, cuantificar prevención efectiva de estafas ni evaluar la precisión del motor de riesgo. Se reconocen limitaciones de representatividad, posibles sesgos de autoselección y ausencia de verificación externa de las experiencias declaradas. El cuestionario se reproduce en el apéndice~\ref{apx:encuesta} y los resultados se exponen en el capítulo~\ref{cap:resultados}. La entrevista con especialista forma parte de las actividades previstas pero no se presenta como realizada.

